\chapter{Turunan} \label{der:chapter}

%%%%%%%%%%%%%%%%%%%%%%%%%%%%%%%%%%%%%%%%%%%%%%%%%%%%%%%%%%%%%%%%%%%%%%%%%%%%%%

\section{Turunan}
\label{sec:der}

%mbxINTROSUBSECTION

\sectionnotes{1 pertemuan}

Gagasan turunan adalah sebagai berikut.
Jika grafik suatu fungsi secara lokal tampak seperti garis lurus,
maka kita dapat membicarakan kemiringan garis tersebut.
Kemiringan itu menyatakan laju perubahan
nilai fungsi pada titik tertentu tersebut.
Tentu saja, pembahasan ini mengecualikan fungsi yang memiliki titik sudut atau
diskontinuitas.
Mari kita nyatakan secara tepat.

\subsection{Definisi dan sifat-sifat dasar}

\begin{defn}
Misalkan $I$ suatu interval, misalkan
$f \colon I \to \R$ suatu fungsi, dan misalkan $c \in I$.  Jika
limit
\begin{equation*}
L \coloneqq \lim_{x \to c} \frac{f(x)-f(c)}{x-c} 
\end{equation*}
ada, kita katakan bahwa $f$
\emph{\myindex{diferensiabel}}\index{fungsi!diferensiabel} di
$c$, kita sebut $L$ sebagai \emph{\myindex{turunan}} $f$ di $c$,
dan kita tulis $f'(c) \coloneqq L$.\glsadd{not:derivative}

\medskip

Jika $f$ diferensiabel di setiap $c \in I$, kita cukup mengatakan bahwa
$f$ \emph{diferensiabel}, dan kemudian kita memperoleh fungsi
$f' \colon I \to \R$.
Turunan kadang-kadang ditulis sebagai $\frac{df}{dx}$ atau
$\frac{d}{dx}\bigl( f(x) \bigr)$.

\medskip

Ekspresi $\frac{f(x)-f(c)}{x-c}$ disebut
\emph{\myindex{hasil bagi selisih}}.
\end{defn}

Interpretasi grafis turunan ditampilkan dalam
\figureref{derivfig}.  Grafik kiri menunjukkan garis yang melalui
$\bigl(c,f(c)\bigr)$
dan $\bigl(x,f(x)\bigr)$ dengan kemiringan
$\frac{f(x)-f(c)}{x-c}$, yaitu
\emph{\myindex{garis sekan}}.  Ketika kita mengambil limit saat $x$ menuju $c$,
kita memperoleh grafik kanan, tempat kita melihat
bahwa turunan fungsi
di titik $c$ merupakan kemiringan garis singgung pada grafik $f$
di titik $\bigl(c,f(c)\bigr)$.

\begin{myfigureht}
\myincludepdft{deriv_derivd}{%
Dua diagram grafik suatu fungsi f.  Di sebelah kiri,
titik c dan x ditandai pada sumbu horizontal.
Titik-titik yang bersesuaian pada grafik f ditandai,
dan sebuah garis (garis sekan) digambar melalui keduanya.  Garis itu diberi label
kemiringan sama dengan f dari x dikurangi f dari c, seluruhnya dibagi
x dikurangi c.
Di sebelah kanan, fungsi yang sama digambar, tetapi hanya titik c
yang ditandai.  Sebuah garis yang menyinggung grafik f
melalui titik yang bersesuaian digambar dan diberi label
kemiringan sama dengan f aksen dari c.}
\caption{Interpretasi grafis turunan.\label{derivfig}}
\end{myfigureht}

Kita memperbolehkan $I$ berupa interval tertutup dan memperbolehkan
$c$ menjadi titik ujung $I$.  Beberapa buku kalkulus tidak memperbolehkan $c$ menjadi
titik ujung interval, tetapi seluruh teorinya tetap berlaku jika hal ini diperbolehkan,
dan pekerjaan kita menjadi lebih mudah.

\begin{example}
Misalkan $f(x) \coloneqq x^2$ didefinisikan pada seluruh garis bilangan real.  Ambil sebarang $c \in \R$.  Kita peroleh bahwa jika
$x \neq c$,
\begin{equation*}
\frac{x^2-c^2}{x-c} =
\frac{(x+c)(x-c)}{x-c} =
(x+c) .
\end{equation*}
Oleh karena itu,
\begin{equation*}
f'(c) = 
\lim_{x\to c} \frac{x^2-c^2}{x-c} =
\lim_{x\to c} (x+c) = 2c.
\end{equation*}
\end{example}

\begin{example}
Misalkan $f(x) \coloneqq ax + b$ untuk bilangan $a, b \in \R$.
Ambil sebarang $c \in \R$.
Untuk $x \neq c$,
\begin{equation*}
\frac{f(x)-f(c)}{x-c} =
\frac{a(x-c)}{x-c} = a .
\end{equation*}
Oleh karena itu,
\begin{equation*}
f'(c) =
\lim_{x\to c} 
\frac{f(x)-f(c)}{x-c} =
\lim_{x\to c} 
a = a.
\end{equation*}
Pada kenyataannya, setiap fungsi diferensiabel secara \myquote{infinitesimal} berperilaku seperti
fungsi afin $ax + b$.  Banyak hasil
dan rumus turunan dapat ditebak dengan terlebih dahulu mengerjakannya untuk fungsi afin.
\end{example}

\begin{example}
Fungsi $f(x) \coloneqq \sqrt{x}$ diferensiabel untuk $x > 0$.  Untuk melihatnya,
tetapkan $c > 0$,
dan andaikan $x \neq c$ serta $x > 0$.  Hitung
\begin{equation*}
\frac{\sqrt{x}-\sqrt{c}}{x-c}
=
\frac{\sqrt{x}-\sqrt{c}}{(\sqrt{x}-\sqrt{c})(\sqrt{x}+\sqrt{c})}
=
\frac{1}{\sqrt{x}+\sqrt{c}} .
\end{equation*}
Oleh karena itu,
\begin{equation*}
f'(c) =
\lim_{x\to c}
\frac{\sqrt{x}-\sqrt{c}}{x-c}
=
\lim_{x\to c}
\frac{1}{\sqrt{x}+\sqrt{c}}
=
\frac{1}{2\sqrt{c}} .
\end{equation*}
\end{example}

\begin{example}
Fungsi $f(x) \coloneqq \sabs{x}$ tidak diferensiabel
di titik asal.  Ketika $x > 0$,
\begin{equation*}
\frac{\sabs{x}-\sabs{0}}{x-0} =
\frac{x-0}{x-0} = 1 .
\end{equation*}
Ketika $x < 0$,
\begin{equation*}
\frac{\sabs{x}-\sabs{0}}{x-0} =
\frac{-x-0}{x-0} = -1 .
\end{equation*}
\end{example}

Sebuah contoh terkenal dari Weierstrass menunjukkan bahwa terdapat fungsi kontinu
yang tidak diferensiabel di \emph{titik mana pun}.  Konstruksi
fungsi tersebut berada di luar cakupan bab ini.  Di sisi lain,
fungsi diferensiabel selalu kontinu.

\begin{prop}
Jika $f \colon I \to \R$ diferensiabel di $c \in I$,
maka fungsi tersebut kontinu di $c$.
\end{prop}

\begin{proof}
Kita mengetahui bahwa limit-limit
\begin{equation*}
\lim_{x\to c}\frac{f(x)-f(c)}{x-c} = f'(c)
\qquad
\text{dan}
\qquad
\lim_{x\to c}(x-c) = 0
\end{equation*}
ada.  Lebih lanjut,
\begin{equation*}
f(x)-f(c) = 
\left( \frac{f(x)-f(c)}{x-c} \right) (x-c) .
\end{equation*}
Oleh karena itu, limit $f(x)-f(c)$ ada dan
\begin{equation*}
\lim_{x\to c} \bigl( f(x)-f(c) \bigr) =
\left(\lim_{x\to c} \frac{f(x)-f(c)}{x-c} \right)
\left(\lim_{x\to c} (x-c) \right) =
f'(c) \cdot 0  = 0.
\end{equation*}
Jadi, $\lim\limits_{x\to c} f(x) = f(c)$, dan $f$ kontinu di $c$.
\end{proof}

Salah satu sifat penting turunan adalah linearitas.  Turunan
merupakan aproksimasi suatu fungsi oleh garis lurus.
Kemiringan garis yang melalui dua titik berubah secara linear ketika
koordinat-$y$ diubah secara linear.  Setelah mengambil limit,
wajar bahwa turunan bersifat linear.

\begin{prop}[Linearitas]
\index{linearitas turunan}
Misalkan $I$ suatu interval, misalkan
$f \colon I \to \R$ dan $g \colon I \to \R$ diferensiabel di $c \in I$,
dan misalkan $\alpha \in \R$.
\begin{enumerate}[(i)]
\item
Definisikan $h \colon I \to \R$ dengan $h(x) \coloneqq \alpha f(x)$.  Maka
$h$ diferensiabel di $c$ dan
$h'(c) = \alpha f'(c)$.
\item
Definisikan $h \colon I \to \R$ dengan $h(x) \coloneqq  f(x) + g(x)$.  Maka
$h$ diferensiabel di $c$ dan
$h'(c) =  f'(c) + g'(c)$.
\end{enumerate}
\end{prop}

\begin{proof}
Pertama, misalkan $h(x) \coloneqq \alpha f(x)$.
Untuk $x \in I$, $x \neq c$,
\begin{equation*}
\frac{h(x)-h(c)}{x-c} =
\frac{\alpha f(x) - \alpha f(c)}{x-c}
=
\alpha \frac{f(x) - f(c)}{x-c} .
\end{equation*}
Limit ruas kanan ketika $x$ menuju $c$ ada
berdasarkan \corref{falg:cor}.  Kita peroleh
\begin{equation*}
\lim_{x\to c}\frac{h(x)-h(c)}{x-c} =
\alpha \lim_{x\to c} \frac{f(x) - f(c)}{x-c} .
\end{equation*}
Oleh karena itu, $h$ diferensiabel di $c$,
dan turunannya dihitung seperti yang dinyatakan.

Berikutnya, definisikan $h(x) \coloneqq f(x)+g(x)$.
Untuk $x \in I$, $x \neq c$, berlaku
\begin{equation*}
\frac{h(x)-h(c)}{x-c} =
\frac{\bigl(f(x) + g(x)\bigr) - \bigl(f(c) + g(c)\bigr)}{x-c}
=
\frac{f(x) - f(c)}{x-c}
+
\frac{g(x) - g(c)}{x-c} .
\end{equation*}
Limit ruas kanan ketika $x$ menuju $c$ ada
berdasarkan \corref{falg:cor}.  Kita peroleh
\begin{equation*}
\lim_{x\to c}\frac{h(x)-h(c)}{x-c} =
\lim_{x\to c} \frac{f(x) - f(c)}{x-c}
+
\lim_{x\to c}\frac{g(x) - g(c)}{x-c} .
\end{equation*}
Oleh karena itu, $h$ diferensiabel di $c$,
dan turunannya dihitung seperti yang dinyatakan.
\end{proof}

Turunan dari hasil kali dua fungsi tidak sama dengan
hasil kali turunannya.
Sebagai gantinya, kita memperoleh apa yang disebut \emph{aturan hasil kali}
atau \emph{\myindex{aturan Leibniz}}%
\footnote{Dinamai menurut matematikawan Jerman
\href{https://en.wikipedia.org/wiki/Leibniz}{Gottfried Wilhelm Leibniz}
(1646--1716).}.

\begin{prop}[Aturan hasil kali]\index{aturan hasil kali}
Misalkan $I$ suatu interval, misalkan
$f \colon I \to \R$ dan $g \colon I \to \R$ merupakan
fungsi yang diferensiabel di $c$.  Jika $h \colon I \to \R$
didefinisikan oleh
\begin{equation*}
h(x) \coloneqq f(x) g(x) ,
\end{equation*}
maka $h$ diferensiabel di $c$ dan
\begin{equation*}
h'(c) = f(c) g'(c) + f'(c) g(c) .
\end{equation*}
\end{prop}

Pembuktian aturan hasil kali diserahkan sebagai latihan.  Kunci pembuktiannya adalah
identitas
$f(x) g(x) - f(c) g(c) =
f(x)\bigl( g(x) - g(c) \bigr)
+ \bigl( f(x) - f(c) \bigr) g(c)$,
yang diilustrasikan dalam \figureref{figprodrule}.
\begin{myfigureht}
\myincludepdft{figprodrule}{%
Diagram suatu daerah persegi panjang pada bidang yang dibagi
menjadi 3 persegi panjang.  Sisi horizontal membentang dari 0 sampai
f dari x, dan sisi vertikal membentang dari 0 sampai g dari x.
Daerah putih terbesar berupa persegi panjang yang membentang secara horizontal
dari 0 sampai f dari c, yang sedikit lebih kecil daripada f dari x,
dan secara vertikal dari 0 sampai g dari c, yang sedikit lebih kecil daripada
g dari x.  Daerah putih ini diberi label f dari c dikali g dari c.
Di atasnya terdapat persegi panjang tipis berarsir terang yang membentang
pada sisi horizontal dari 0 sampai f dari x dan pada sisi vertikal
dari g dari c sampai g dari x.  Daerah itu diberi label
f dari x dikali besaran g dari x dikurangi g dari c.
Terakhir, terdapat persegi panjang tipis berarsir gelap
yang membentang secara horizontal dari f dari c sampai f dari x dan secara vertikal
dari 0 sampai g dari c.  Daerah itu diberi label f dari x dikurangi f dari c,
seluruhnya dikali g dari c.}
\caption{Gagasan aturan hasil kali.  Luas seluruh persegi panjang
$f(x)g(x)$ berbeda dari luas persegi panjang putih $f(c)g(c)$
sebesar luas persegi panjang berarsir terang
$f(x)\bigl( g(x) - g(c) \bigr)$ ditambah persegi panjang yang lebih gelap
$\bigl( f(x) - f(c) \bigr) g(c)$.
Dengan kata lain, secara kasar, $\Delta (f \cdot g)
= f \cdot \Delta g + \Delta f \cdot g$.\label{figprodrule}}
\end{myfigureht}



\begin{prop}[Aturan hasil bagi]\index{aturan hasil bagi}
Misalkan $I$ suatu interval, misalkan
$f \colon I \to \R$ dan $g \colon I \to \R$ diferensiabel di $c$
dan $g(x) \neq 0$ untuk semua $x \in I$.
Jika $h \colon I \to \R$
didefinisikan oleh
\begin{equation*}
h(x) \coloneqq \frac{f(x)}{g(x)},
\end{equation*}
maka $h$ diferensiabel di $c$ dan
\begin{equation*}
h'(c) = \frac{f'(c) g(c) - f(c) g'(c)}{{\bigl(g(c)\bigr)}^2} .
\end{equation*}
\end{prop}

Sekali lagi, pembuktiannya diserahkan sebagai latihan.

\subsection{Aturan rantai}

Fungsi yang lebih rumit sering kali diperoleh melalui komposisi,
yang diturunkan dengan menggunakan aturan rantai.  Aturan ini juga memberi tahu kita
bagaimana suatu turunan berubah jika kita melakukan perubahan variabel.

\begin{prop}[Aturan rantai]
\index{aturan rantai}
Misalkan $I_1, I_2$ adalah interval, dan misalkan
$g \colon I_1 \to I_2$ diferensiabel di $c \in I_1$,
serta
$f \colon I_2 \to \R$ diferensiabel di $g(c)$.
Jika $h \colon I_1 \to \R$
didefinisikan oleh
\begin{equation*}
h(x) \coloneqq (f \circ g) (x) = f\bigl(g(x)\bigr) ,
\end{equation*}
maka $h$ diferensiabel di $c$ dan
\begin{equation*}
h'(c) = f'\bigl(g(c)\bigr)g'(c) .
\end{equation*}
\end{prop}

\begin{proof}
Misalkan $d \coloneqq g(c)$.  Definisikan
$u \colon I_2 \to \R$ dan $v \colon I_1 \to \R$ dengan
\begin{equation*}
u(y) \coloneqq
\begin{cases}
 \frac{f(y) - f(d)}{y-d}  & \text{jika } y \neq d, \\
 f'(d)                    & \text{jika } y = d,
\end{cases}
\qquad
v(x) \coloneqq
\begin{cases}
 \frac{g(x) - g(c)}{x-c} & \text{jika } x \neq c, \\
 g'(c)                   & \text{jika } x = c.
\end{cases}
\end{equation*}
Karena $f$ diferensiabel di $d = g(c)$, kita peroleh bahwa
$u$ kontinu di $d$.  Demikian pula, $v$ kontinu di $c$.
Untuk setiap $x$ dan $y$,
\begin{equation*}
f(y)-f(d) = u(y) (y-d)
\qquad \text{dan} \qquad
g(x)-g(c) = v(x) (x-c) .
\end{equation*}
Substitusikan persamaan-persamaan ini untuk memperoleh
\begin{equation*}
h(x)-h(c)
=
f\bigl(g(x)\bigr)-f\bigl(g(c)\bigr)
=
u\bigl( g(x) \bigr) \bigl(g(x)-g(c)\bigr)
=
u\bigl( g(x) \bigr) \bigl(v(x) (x-c)\bigr) .
\end{equation*}
Oleh karena itu, jika $x \neq c$,
\begin{equation} \label{eq:chainruleeq}
\frac{h(x)-h(c)}{x-c}
=
u\bigl( g(x) \bigr) v(x) .
\end{equation}
Dari kekontinuan $u$ dan $v$ masing-masing di $d$ dan $c$, kita peroleh
$\lim_{y \to d} u(y)
= f'(d) = f'\bigl(g(c)\bigr)$ dan
$\lim_{x \to c} v(x) = g'(c)$.
Fungsi $g$ kontinu di $c$, sehingga $\lim_{x \to c} g(x) = g(c)$.
Jadi, limit
ruas kanan \eqref{eq:chainruleeq}
saat $x$ menuju $c$
ada dan sama dengan $f'\bigl(g(c)\bigr) g'(c)$.  Dengan demikian, $h$
diferensiabel di $c$ dan $h'(c) = f'\bigl(g(c)\bigr)g'(c)$.
\end{proof}

\subsection{Latihan}

\begin{exercise}
Buktikan aturan hasil kali.
Petunjuk: Buktikan dan gunakan
$f(x) g(x) - f(c) g(c) = f(x)\bigl( g(x) - g(c) \bigr) + \bigl( f(x) -
f(c) \bigr) g(c)$.
\end{exercise}

\begin{exercise}
Buktikan aturan hasil bagi.  Petunjuk: Anda dapat melakukannya secara langsung, tetapi mungkin
lebih mudah untuk mencari turunan $\nicefrac{1}{x}$ lalu menggunakan
aturan rantai dan aturan hasil kali.
\end{exercise}

\begin{exercise} \label{exercise:diffofxn}
Untuk $n \in \Z$,
buktikan bahwa $x^n$ diferensiabel dan carilah turunannya,
kecuali, tentu saja, jika $n < 0$ dan $x=0$.
Petunjuk: Gunakan aturan hasil kali.
\end{exercise}

\begin{exercise}
Buktikan bahwa suatu polinomial diferensiabel, dan carilah turunannya.
Petunjuk: Gunakan latihan sebelumnya.
\end{exercise}

\begin{exercise}
Definisikan $f \colon \R \to \R$ dengan
\begin{equation*}
f(x) \coloneqq
\begin{cases}
x^2 & \text{jika } x \in \Q,\\
0 & \text{selainnya.}
\end{cases}
\end{equation*}
Buktikan bahwa $f$ diferensiabel di $0$, tetapi diskontinu di semua titik
kecuali $0$.
\end{exercise}

\begin{exercise}
Asumsikan ketaksamaan $\babs{x-\sin(x)} \leq x^2$.  Buktikan bahwa $\sin$
diferensiabel di $0$, dan carilah turunannya di $0$.
\end{exercise}

\begin{exercise}
Dengan menggunakan latihan sebelumnya, buktikan bahwa $\sin$ diferensiabel di setiap $x$
dan bahwa turunannya adalah $\cos(x)$.  Petunjuk: Gunakan identitas trigonometri
pengubahan jumlah menjadi hasil kali sebagaimana yang kita lakukan sebelumnya.
\end{exercise}

\begin{exercise}
Misalkan $f \colon I \to \R$ diferensiabel.  Untuk $n \in \Z$, misalkan $f^n$
adalah fungsi yang didefinisikan oleh $f^n(x) \coloneqq {\bigl( f(x) \bigr)}^n$.  Jika
$n < 0$, asumsikan $f(x) \neq 0$ untuk setiap $x \in I$.  Buktikan bahwa
$(f^n)'(x) = n {\bigl(f(x) \bigr)}^{n-1} f'(x)$.
\end{exercise}

\begin{exercise}
Misalkan $f \colon \R \to \R$ adalah fungsi kontinu Lipschitz yang
diferensiabel.
Buktikan bahwa $f'$ adalah fungsi terbatas.
\end{exercise}

\begin{exercise} \label{exercise:inversederformula}
Misalkan $I_1, I_2$ adalah interval.
Misalkan $f \colon I_1 \to I_2$ adalah fungsi bijektif dan $g \colon I_2 \to I_1$
adalah inversnya.  Misalkan $f$ diferensiabel di $c \in I_1$,
$f'(c) \neq 0$, dan $g$ diferensiabel di $f(c)$.  Gunakan aturan rantai
untuk mencari rumus bagi $g'\bigl(f(c)\bigr)$ (dinyatakan dalam $f'(c)$).
\end{exercise}

\begin{exercise} \label{exercise:bndmuldiff}
Misalkan $f \colon I \to \R$ terbatas, $g \colon I \to
\R$ diferensiabel di $c \in I$, dan $g(c) = g'(c) = 0$.  Tunjukkan
bahwa $h(x) \coloneqq f(x) g(x)$ diferensiabel di $c$.  Petunjuk: Anda
tidak dapat menerapkan aturan hasil kali.
\end{exercise}

\begin{exercise} \label{exercise:diffsqueeze}
Misalkan $f \colon I \to \R$, 
$g \colon I \to \R$, dan
$h \colon I \to \R$ adalah fungsi-fungsi.  Misalkan $c \in I$ sedemikian sehingga
$f(c) = g(c) = h(c)$, $g$ dan $h$ diferensiabel di $c$,
dan $g'(c) = h'(c)$.  Selanjutnya, misalkan $h(x) \leq f(x) \leq g(x)$ untuk
setiap $x \in I$.  Buktikan bahwa $f$ diferensiabel di $c$ dan $f'(c) = g'(c) =
h'(c)$.
\end{exercise}

\begin{exercise}
\leavevmode
\begin{enumerate}[a)]
\item
Misalkan $f \colon (-1,1) \to \R$ adalah fungsi sedemikian sehingga $f(x) = x h(x)$ untuk
suatu fungsi terbatas $h \colon (-1,1) \to \R$.
Tunjukkan bahwa $g(x) \coloneqq {\bigl( f(x) \bigr)}^2$
diferensiabel di titik asal dan $g'(0) = 0$.
\item
Carilah contoh suatu
fungsi kontinu $f \colon (-1,1) \to \R$ dengan $f(0) = 0$, tetapi sedemikian
sehingga $g(x) \coloneqq {\bigl( f(x) \bigr)}^2$ tidak diferensiabel di titik asal.
\end{enumerate}
\end{exercise}

\begin{exercise}
Misalkan $f \colon I \to \R$ diferensiabel di $c \in I$.
Buktikan bahwa terdapat bilangan $a$ dan $b$ dengan sifat bahwa
untuk setiap $\epsilon > 0$, terdapat $\delta > 0$ sedemikian sehingga
$\babs{a+b(x-c) - f(x)} \leq \epsilon \sabs{x-c}$, apabila $x \in I$ dan
$\sabs{x-c} < \delta$.
Dengan kata lain, tunjukkan bahwa
terdapat fungsi $g \colon I \to \R$
sedemikian sehingga $\lim_{x\to c} g(x) = 0$ dan
$\babs{a+b(x-c) - f(x)} = g(x) \sabs{x-c}$.
\end{exercise}

\begin{exercise} \label{exercise:simpleLHopital}
Buktikan versi sederhana dari \myindex{aturan L'H\^opital}\index{aturan L'Hospital} berikut.
Misalkan 
$f \colon (a,b) \to \R$ dan $g \colon (a,b) \to \R$ adalah fungsi-fungsi
diferensiabel
yang turunannya, $f'$ dan $g'$, merupakan fungsi kontinu.
Misalkan di $c \in (a,b)$ berlaku $f(c) = 0$, $g(c)=0$,
$g'(x) \neq 0$ untuk setiap $x \in (a,b)$, dan
$g(x) \neq 0$ apabila $x \neq c$.   Perhatikan
bahwa limit $\nicefrac{f'(x)}{g'(x)}$ saat $x$ menuju $c$ ada.  Tunjukkan bahwa
\begin{equation*}
\lim_{x \to c} \frac{f(x)}{g(x)} = 
\lim_{x \to c} \frac{f'(x)}{g'(x)} .
\end{equation*}
\end{exercise}

\begin{exercise}
Misalkan $f \colon (a,b) \to \R$ diferensiabel di $c \in (a,b)$,
$f(c)=0$, dan $f'(c) > 0$.
Buktikan bahwa terdapat $\delta > 0$ sedemikian sehingga $f(x) < 0$
apabila $c-\delta < x < c$ dan
$f(x) > 0$ 
apabila $c < x < c+\delta$.
\end{exercise}

%%%%%%%%%%%%%%%%%%%%%%%%%%%%%%%%%%%%%%%%%%%%%%%%%%%%%%%%%%%%%%%%%%%%%%%%%%%%%%

\sectionnewpage
\section{Teorema nilai rata-rata}
\label{sec:mvt}

%mbxINTROSUBSECTION

\sectionnotes{2 kuliah (beberapa penerapan dapat dilewati)}

\subsection{Minimum dan maksimum relatif}

Sebelumnya kita telah membahas maksimum dan minimum mutlak.  Keduanya merupakan puncak tertinggi dan
lembah terendah di seluruh jajaran pegunungan.  Bagaimana dengan
puncak masing-masing gunung dan dasar masing-masing lembah?
Turunan, sebagai konsep lokal, ibarat berjalan di tengah kabut; turunan
tidak dapat memberi tahu apakah kita berada di puncak tertinggi, tetapi dapat memberi tahu apakah kita berada
di puncak suatu gunung.

\begin{defn}
Misalkan $S \subset \R$ adalah suatu himpunan dan
misalkan $f \colon S \to \R$ adalah suatu fungsi.  Fungsi $f$ dikatakan mempunyai
\emph{\myindex{maksimum relatif}}\index{maksimum!relatif}
di $c \in S$ jika terdapat $\delta>0$
sedemikian sehingga untuk setiap $x \in S$ dengan $\sabs{x-c} < \delta$,
berlaku $f(x) \leq f(c)$.
Definisi
\emph{\myindex{minimum relatif}}\index{minimum!relatif}
serupa.
\end{defn}

\begin{lemma}\label{relminmax:lemma}
Misalkan $f \colon (a,b) \to \R$ diferensiabel di $c \in (a,b)$,
dan $f$ mempunyai
minimum relatif atau maksimum relatif di $c$.  Maka
$f'(c) = 0$.
\end{lemma}

\begin{proof}
Misalkan $c$ adalah 
titik maksimum relatif bagi $f$.  Artinya, terdapat $\delta > 0$ sedemikian
sehingga untuk setiap $x \in (a,b)$ dengan
$\sabs{x-c} < \delta$, berlaku $f(x)-f(c) \leq 0$.
Perhatikan hasil bagi
selisih.  Jika $c < x < c+\delta$, maka
\begin{equation*}
\frac{f(x)-f(c)}{x-c} \leq 0 ,
\end{equation*}
dan jika $c-\delta < y < c$, maka
\begin{equation*}
\frac{f(y)-f(c)}{y-c} \geq 0 .
\end{equation*}
Lihat \figureref{fig:critpt} untuk ilustrasinya.
\begin{myfigureht}
\myincludegraphics{critpt}{%
Ditampilkan grafik suatu fungsi f.
Tiga titik ditandai pada sumbu horizontal, dari kiri
ke kanan, yaitu y, c, dan x.  Fungsi tersebut mempunyai maksimum
di titik c.  Titik-titik pada grafik yang bersesuaian dengan x, c, dan y
ditandai.
Melalui titik-titik yang bersesuaian
dengan y dan c terdapat sebuah garis sekan yang menanjak dan
diberi keterangan bahwa kemiringannya sama dengan f dari y dikurangi f dari c, seluruhnya
dibagi besaran y-c, lebih besar daripada atau sama dengan 0.
Melalui titik-titik yang bersesuaian
dengan c dan x terdapat sebuah garis sekan yang menurun dan
diberi keterangan bahwa kemiringannya sama dengan f dari x dikurangi f dari c, seluruhnya
dibagi besaran x-c, lebih kecil daripada atau sama dengan 0.}
\caption{Kemiringan garis-garis sekan pada suatu maksimum relatif.\label{fig:critpt}}
\end{myfigureht}

Karena $a < c < b$, terdapat
barisan $\{ x_n \}_{n=1}^\infty$ dan
$\{ y_n \}_{n=1}^\infty$ di $(a,b)$
sedemikian sehingga $c < x_n < c+\delta$ dan
$c-\delta < y_n < c$ untuk setiap $n \in \N$, serta
$\lim_{n\to\infty} x_n = \lim_{n\to\infty} y_n = c$.
Karena $f$
diferensiabel di $c$,
\begin{equation*}
0 \geq \lim_{n\to\infty} \frac{f(x_n)-f(c)}{x_n-c} 
=
f'(c)
=
\lim_{n\to\infty} \frac{f(y_n)-f(c)}{y_n-c} \geq 0.
\end{equation*}
Bukti untuk kasus maksimum selesai.
Untuk kasus minimum, perhatikan
fungsi $-f$.
\end{proof}

Untuk suatu fungsi diferensiabel, titik tempat 
$f'(c) = 0$ disebut \emph{\myindex{titik kritis}}.  Ketika $f$ tidak
diferensiabel di beberapa titik,
lazim pula dikatakan bahwa $c$ adalah titik kritis
jika $f'(c)$ tidak ada.
Teorema tersebut menyatakan bahwa minimum atau maksimum relatif di titik interior
suatu interval haruslah merupakan titik kritis.
Seperti yang Anda ingat dari kalkulus, minimum dan maksimum suatu fungsi dicari
dengan mencari semua titik kritis beserta titik-titik ujung
interval tersebut, lalu memeriksa di titik mana di antara titik-titik itu
fungsi bernilai terbesar atau terkecil.

\subsection{Teorema Rolle}

Andaikan suatu fungsi memiliki nilai yang sama di kedua titik ujung suatu interval.
Secara intuitif, fungsi tersebut semestinya mencapai minimum atau maksimum di bagian dalam
interval.
Pada minimum atau maksimum tersebut, turunannya semestinya bernilai nol.
Lihat \figureref{rollefig} untuk gagasan geometrisnya.
Inilah isi teorema Rolle%
\footnote{Dinamai menurut matematikawan Prancis
\href{https://en.wikipedia.org/wiki/Michel_Rolle}{Michel Rolle}
(1652--1719).}.

\begin{myfigureht}
\myincludegraphics{rollefig}{%
Grafik suatu fungsi digambar pada interval dari a sampai b.
Fungsi bernilai nol di a dan b.  Fungsi naik menuju
suatu maksimum,
lalu turun menuju suatu minimum, kemudian naik kembali.  Maksimum fungsi terjadi
di titik yang ditandai c.  Garis singgung horizontal digambar pada
titik yang bersesuaian di grafik.  Perhatikan bahwa garis itu sejajar dengan
sumbu horizontal.}
\caption{Titik tempat garis singgung horizontal, yaitu $f'(c) =
0$.\label{rollefig}}
\end{myfigureht}

\begin{thm}[Rolle] \label{thm:rolle}
\index{teorema Rolle}
Misalkan $f \colon [a,b] \to \R$ suatu fungsi kontinu
yang diferensiabel pada $(a,b)$ sedemikian sehingga $f(a) = f(b)$.
Maka terdapat $c \in (a,b)$ sedemikian sehingga $f'(c) = 0$.
\end{thm}

\begin{proof}
Karena $f$ kontinu pada $[a,b]$, fungsi tersebut mencapai minimum mutlak dan
maksimum mutlak pada $[a,b]$.  Kita ingin menerapkan \lemmaref{relminmax:lemma},
sehingga kita perlu menemukan $c \in (a,b)$ tempat $f$ mencapai minimum atau
maksimum.
Tuliskan $K \coloneqq f(a) = f(b)$.
Jika terdapat $x$ sedemikian sehingga $f(x) > K$, maka maksimum mutlak
lebih besar daripada $K$ dan karenanya terjadi di suatu $c \in (a,b)$,
sehingga $f'(c) = 0$.  Di sisi lain, jika terdapat $x$
sedemikian sehingga $f(x) < K$, maka minimum mutlak terjadi di suatu
$c \in (a,b)$, sehingga $f'(c) = 0$.  Jika tidak ada $x$ sedemikian sehingga
$f(x) > K$ atau
$f(x) < K$, maka $f(x) = K$ untuk semua $x$, dan dengan demikian
$f'(x) = 0$ untuk semua $x \in (a,b)$, sehingga sebarang $c \in (a,b)$ memenuhi.
\end{proof}

Syarat bahwa turunan ada untuk semua $x
\in (a,b)$ mutlak diperlukan.  Perhatikan fungsi $f(x) \coloneqq \sabs{x}$ pada $[-1,1]$.
Jelas bahwa $f(-1) = f(1)$, tetapi tidak ada titik $c$ tempat $f'(c) = 0$.

\subsection{Teorema nilai rata-rata}

Kita memperluas \hyperref[thm:rolle]{teorema Rolle}
ke fungsi-fungsi yang mencapai nilai yang berbeda
di kedua titik ujungnya.

\begin{thm}[Teorema nilai rata-rata] \label{thm:mvt}
\index{teorema nilai rata-rata}
Misalkan $f \colon [a,b] \to \R$ merupakan fungsi kontinu yang
diferensiabel pada $(a,b)$.  Maka terdapat titik $c \in (a,b)$
sedemikian sehingga
\begin{equation*}
f(b)-f(a) = f'(c)(b-a) .
\end{equation*}
\end{thm}

Untuk interpretasi geometris teorema nilai rata-rata, lihat
\figureref{mvtfig}.  Gagasannya ialah bahwa nilai $\frac{f(b)-f(a)}{b-a}$
merupakan kemiringan garis yang melalui titik-titik $\bigl(a,f(a)\bigr)$
dan $\bigl(b,f(b)\bigr)$.
Kemudian $c$ adalah titik sedemikian sehingga $f'(c) = \frac{f(b)-f(a)}{b-a}$, yaitu 
garis singgung di titik $\bigl(c,f(c)\bigr)$ memiliki kemiringan yang sama dengan
garis yang melalui $\bigl(a,f(a)\bigr)$ dan $\bigl(b,f(b)\bigr)$.
Nama tersebut berasal dari fakta bahwa kemiringan garis sekan
adalah nilai rata-rata turunan, sehingga turunan mencapai nilai rata-ratanya
di bagian dalam interval.

Teorema ini mengikuti dari \hyperref[thm:rolle]{teorema Rolle}
dengan mengurangkan dari $f$ fungsi linear afin
yang nilainya di $a$ dan $b$ sama dengan nilai $f$.
Grafik fungsi linear afin ini adalah 
garis sekan---garis lurus yang melalui
$\bigl(a,f(a)\bigr)$ dan $\bigl(b,f(b)\bigr)$---dan
turunannya adalah
$\frac{f(b)-f(a)}{b-a}$.
Dengan demikian, kita mencari titik tempat
selisih ini memiliki turunan nol.

\begin{myfigureht}
\myincludegraphics{mvtfig}{%
 Grafik suatu fungsi digambar pada interval dari a hingga b.
 Fungsi tersebut tampak serupa dengan sebelumnya, tetapi tidak lagi
 bernilai nol di titik-titik ujung; titik ujung kiri bernilai negatif
 dan titik ujung kanan bernilai positif.  Digambar sebuah garis lurus yang melalui kedua titik ujung
 kurva, yakni garis yang melalui titik
 (a, f dari a) dan (b, f dari b).  Di titik yang bersesuaian dengan 
 c pada sumbu horizontal, garis singgung pada fungsi
 sejajar dengan garis yang melalui kedua titik ujung.}
\caption{Interpretasi grafis teorema nilai rata-rata.\label{mvtfig}}
\end{myfigureht}


\begin{proof}
Definisikan
fungsi $g \colon [a,b] \to \R$ dengan
\begin{equation*}
g(x) \coloneqq
f(x)-
\left( f(b)+\frac{f(b)-f(a)}{b-a}(x-b) \right)
=
f(x)-
f(b)-\frac{f(b)-f(a)}{b-a}(x-b) .
\end{equation*}
Fungsi $g$ diferensiabel pada $(a,b)$,
kontinu pada $[a,b]$, serta $g(a) = 0$ dan $g(b) = 0$.  Jadi, terdapat
suatu
$c \in (a,b)$ sedemikian sehingga $g'(c) = 0$, yaitu,
\begin{equation*}
0 = g'(c) = f'(c)-\frac{f(b)-f(a)}{b-a} .
\end{equation*}
Dengan kata lain,
$f(b)-f(a) = f'(c)(b-a)$.
\end{proof}

Pembuktian tersebut dapat diperumum.  Dengan meninjau
$g(x) \coloneqq
\bigl(f(x)-f(b)\bigr)
\bigl(\varphi(b)-\varphi(a)\bigr)
-
\bigl(f(b)-f(a)\bigr)
\bigl(\varphi(x)-\varphi(b)\bigr)$,
dapat dibuktikan versi berikut.  Pembuktiannya kita tinggalkan sebagai latihan.

\begin{thm}[Teorema nilai rata-rata Cauchy] \label{thm:cauchymvt}
\index{teorema nilai rata-rata Cauchy}
Misalkan $f \colon [a,b] \to \R$ dan $\varphi \colon [a,b] \to \R$ merupakan
fungsi-fungsi kontinu yang
diferensiabel pada $(a,b)$.  Maka terdapat titik $c \in (a,b)$
sedemikian sehingga
\begin{equation*}
\bigl(f(b)-f(a)\bigr)\varphi'(c) = f'(c)\bigl(\varphi(b)-\varphi(a)\bigr) .
\end{equation*}
\end{thm}

Teorema nilai rata-rata memiliki keistimewaan sebagai salah satu dari sedikit teorema yang
pernah dikutip
di pengadilan.  Ketika polisi mengukur kecepatan mobil dari pesawat terbang, atau
melalui kamera yang membaca pelat nomor, mereka 
mengukur waktu yang dibutuhkan mobil untuk bergerak di antara dua titik.
Teorema nilai rata-rata kemudian
menyatakan bahwa di suatu tempat mobil itu pasti pernah mencapai kecepatan yang diperoleh dari hasil bagi
selisih jarak dan selisih waktu.



\subsection{Penerapan}

Mari kita lihat beberapa penerapan teorema nilai rata-rata.
Penerapan-penerapan ini memperlihatkan penggunaan khas teorema tersebut, yaitu
menghilangkan suatu limit dengan menemukan jenis titik yang tepat, tempat
turunannya bukan sekadar mendekati suatu hasil bagi selisih, melainkan benar-benar
sama dengannya.
Pertama, kita menyelesaikan persamaan diferensial pertama kita.

\begin{prop} \label{prop:derzeroconst}
Misalkan $I$ suatu interval dan
$f \colon I \to \R$ suatu fungsi diferensiabel sedemikian sehingga $f'(x) = 0$
untuk semua $x \in I$.
Maka $f$ konstan.
\end{prop}

\begin{proof}
Ambil sebarang $x,y \in I$ dengan $x < y$.
Karena $I$ suatu interval, $[x,y] \subset I$.
Maka pembatasan $f$ pada $[x,y]$ memenuhi hipotesis
\hyperref[thm:mvt]{teorema nilai rata-rata}.
Oleh karena itu, terdapat $c \in (x,y)$ sedemikian sehingga
\begin{equation*}
f(y)-f(x) = f'(c)(y-x).
\end{equation*}
Karena $f'(c) = 0$, kita peroleh $f(y) = f(x)$.  Jadi,
fungsi tersebut konstan.
\end{proof}

Setelah mengetahui apa artinya suatu fungsi tetap konstan, kita meninjau
fungsi yang monoton naik dan monoton turun.
Kita katakan bahwa $f \colon I \to \R$ \emph{\myindex{monoton naik}}
(berturut-turut, \emph{\myindex{naik tegas}}) jika
$x < y$ mengakibatkan $f(x) \leq f(y)$ (berturut-turut, $f(x) < f(y)$).
Kita mendefinisikan
\emph{\myindex{monoton turun}} dan
\emph{\myindex{turun tegas}} dengan cara yang sama dengan membalik
pertidaksamaan untuk $f$.

\begin{prop} \label{incdecdiffprop}
Misalkan $I$ suatu interval dan
$f \colon I \to \R$ suatu fungsi diferensiabel.
%\begin{enumerate}[(i),itemsep=0.5\itemsep,parsep=0.5\parsep,topsep=0.5\topsep,partopsep=0.5\partopsep]
\begin{enumerate}[(i)]
\item $f$ monoton naik jika dan hanya jika $f'(x) \geq 0$ untuk semua $x \in I$.
\item $f$ monoton turun jika dan hanya jika $f'(x) \leq 0$ untuk semua $x \in I$.
\end{enumerate}
\end{prop}

\begin{proof}
Mari kita buktikan butir pertama.  Andaikan $f$ monoton naik.
Untuk semua $x,c \in I$ dengan $x \neq c$,
\begin{equation*}
\frac{f(x)-f(c)}{x-c} \geq 0 .
\end{equation*}
Dengan mengambil limit saat $x$ menuju $c$, kita peroleh $f'(c) \geq 0$.

Untuk arah sebaliknya, andaikan $f'(x) \geq 0$ untuk semua $x \in I$.
Ambil sebarang $x, y \in I$ dengan $x < y$, dan perhatikan bahwa $[x,y] \subset I$.
Menurut \hyperref[thm:mvt]{teorema nilai rata-rata}, terdapat $c \in (x,y)$ sedemikian sehingga
\begin{equation*}
f(y)-f(x) = f'(c)(y-x) .
\end{equation*}
Karena $f'(c) \geq 0$ dan $y-x > 0$, kita peroleh $f(y) - f(x) \geq 0$ atau $f(x) \leq
f(y)$, sehingga
$f$ monoton naik.

Butir kedua, yaitu $f$ monoton turun, kita serahkan kepada pembaca sebagai latihan.
\end{proof}

Pernyataan serupa tetapi lebih lemah berlaku untuk fungsi naik tegas dan
fungsi turun tegas.

\begin{prop} \label{incdecdiffstrictprop}
Misalkan $I$ suatu interval dan
$f \colon I \to \R$ suatu fungsi diferensiabel.
\begin{enumerate}[(i)]
\item
\label{incdecdiffstrictprop:i}
Jika $f'(x) > 0$ untuk semua $x \in I$, maka
$f$ naik tegas.
\item
\label{incdecdiffstrictprop:ii}
Jika $f'(x) < 0$ untuk semua $x \in I$,
maka $f$ turun tegas.
\end{enumerate}
\end{prop}

Pembuktian
\ref{incdecdiffstrictprop:i}
kita serahkan sebagai latihan.
Selanjutnya, \ref{incdecdiffstrictprop:ii}
mengikuti dari 
\ref{incdecdiffstrictprop:i} dengan meninjau $-f$
sebagai gantinya.
Kebalikan proposisi ini tidak benar.  Fungsi
$f(x) \coloneqq x^3$ naik tegas, tetapi $f'(0) = 0$.

\medskip

Penerapan lain dari \hyperref[thm:mvt]{teorema nilai rata-rata} adalah hasil berikut mengenai
letak ekstremum, yang kadang disebut
\emph{\myindex{uji turunan pertama}}.
Hasil ini dinyatakan untuk minimum dan maksimum mutlak.
Untuk menerapkannya guna mencari minimum
dan maksimum relatif, batasi $f$ pada interval $(c-\delta,c+\delta)$.

\begin{prop} \label{firstderminmaxtest}
Misalkan $f \colon (a,b) \to \R$ kontinu.  Misalkan $c \in (a,b)$
dan andaikan
$f$ diferensiabel pada $(a,c)$ dan $(c,b)$.
\begin{enumerate}[(i)]
\item Jika $f'(x) \leq 0$ untuk setiap $x \in (a,c)$ dan
 $f'(x) \geq 0$ untuk setiap $x \in (c,b)$, maka $f$ memiliki minimum mutlak 
di $c$.
\item Jika $f'(x) \geq 0$ untuk setiap $x \in (a,c)$ dan
 $f'(x) \leq 0$ untuk setiap $x \in (c,b)$, maka $f$ memiliki maksimum mutlak
di $c$.
\end{enumerate}
\end{prop}

\begin{proof}
Kita membuktikan butir pertama dan menyerahkan butir kedua kepada pembaca.
Ambil $x \in (a,c)$
dan suatu barisan $\{ y_n\}_{n=1}^\infty$ sedemikian sehingga $x < y_n < c$ untuk semua $n$
dan $\lim_{n\to\infty} y_n = c$.
Menurut proposisi sebelumnya,
$f$ monoton turun pada $(a,c)$ sehingga $f(x) \geq f(y_n)$ untuk semua $n$.
Karena $f$
kontinu di $c$, dengan mengambil limit kita peroleh
$f(x) \geq f(c)$.

Serupa dengan itu, ambil $x \in (c,b)$
dan suatu barisan $\{ y_n\}_{n=1}^\infty$ sedemikian sehingga $c < y_n < x$ dan
$\lim_{n\to\infty} y_n = c$.
Fungsi tersebut monoton naik pada $(c,b)$ sehingga $f(x) \geq f(y_n)$ untuk semua $n$.
Dari kekontinuan $f$, kita peroleh
$f(x) \geq f(c)$.  Jadi $f(x) \geq f(c)$ untuk semua
$x \in (a,b)$.
\end{proof}

Kebalikan proposisi tersebut tidak berlaku.  Lihat
\exampleref{baddifffunc:example} di bawah.

\medskip

Penerapan teorema nilai rata-rata lain yang sering digunakan dan mungkin pernah Anda
lihat dalam kalkulus adalah hasil berikut mengenai keterdiferensialan di
titik-titik ujung suatu interval.  Pembuktiannya diberikan dalam
\exerciseref{exercise:endpointderivative}.

\begin{prop} \label{prop:endpointderivative}
\leavevmode
\begin{enumerate}[(i)]
\item
Andaikan $f \colon [a,b) \to \R$ kontinu, diferensiabel pada $(a,b)$,
dan $\lim_{x \to a} f'(x) = L$.  Maka $f$ diferensiabel di $a$ dan
$f'(a) = L$.
\item
Andaikan $f \colon (a,b] \to \R$ kontinu, diferensiabel pada $(a,b)$,
dan $\lim_{x \to b} f'(x) = L$.  Maka $f$ diferensiabel di $b$ dan
$f'(b) = L$.
\end{enumerate}
\end{prop}

Sesungguhnya, dengan menggunakan hasil perluasan \propref{context:prop}, kita tidak perlu mengandaikan
bahwa $f$ terdefinisi di titik ujung.  Lihat \exerciseref{exercise:extendboundedder}.

\subsection{Kekontinuan turunan dan teorema nilai antara}

Turunan fungsi memenuhi
sifat nilai antara.

\begin{thm}[Darboux] \label{thm:darboux} \index{teorema Darboux}
Misalkan $f \colon [a,b] \to \R$ diferensiabel.  Andaikan $y \in \R$ sedemikian
sehingga $f'(a) < y < f'(b)$ atau
$f'(a) > y > f'(b)$.  Maka terdapat $c \in (a,b)$ sedemikian sehingga $f'(c) =
y$.
\end{thm}

Pembuktian diperoleh dengan menjadikan $g$ sebagai selisih $f$
dan suatu fungsi linear dengan turunan $y$.
Fungsi baru $g$ mereduksi masalah
ke kasus $y=0$, dengan $g'(a) > 0 > g'(b)$.  Artinya, $g$ naik di $a$ dan
turun di $b$, sehingga harus mencapai maksimum di dalam $(a,b)$,
tempat turunannya bernilai nol.  Lihat \figureref{darbouxthmfig}.

\begin{myfigureht}
\myincludegraphics{darbouxthmfig}{%
Grafik suatu fungsi dengan tiga titik, a, c, b, dalam urutan tersebut,
yang ditandai pada sumbu horizontal.  Garis singgung digambar pada
titik-titik yang bersesuaian di grafik.  Dari kiri ke kanan,
garis pertama miring naik dan diberi label g prima di a lebih besar daripada 0.
Garis kedua horizontal dan diberi label g prima di c sama dengan 0.
Garis ketiga miring turun dan diberi label g prima di b lebih kecil daripada 0.}
\caption{Gagasan pembuktian teorema Darboux.\label{darbouxthmfig}}
\end{myfigureht}

\begin{proof}
Andaikan
$f'(a) < y < f'(b)$.
Definisikan
\begin{equation*}
g(x) \coloneqq yx - f(x) .
\end{equation*}
Fungsi $g$ kontinu pada $[a,b]$, sehingga $g$ mencapai maksimum di suatu $c \in
[a,b]$.

Fungsi $g$ juga diferensiabel pada $[a,b]$.
Hitung $g'(x) = y-f'(x)$.  Jadi, $g'(a) > 0$.  Karena turunan merupakan
limit hasil bagi selisih dan bernilai positif, pasti terdapat suatu
hasil bagi selisih yang positif.  Artinya, pasti terdapat
$x > a$ sedemikian sehingga
\begin{equation*}
\frac{g(x)-g(a)}{x-a} > 0 ,
\end{equation*}
atau $g(x) > g(a)$.  Jadi, $g$
tidak mungkin memiliki maksimum di $a$.  Demikian pula, karena $g'(b) < 0$,
kita temukan $x < b$ (suatu $x$ yang berbeda) sedemikian sehingga
$\frac{g(x)-g(b)}{x-b} < 0$ atau $g(x) > g(b)$, sehingga
$g$ tidak mungkin memiliki maksimum di $b$.
Oleh karena itu, $c \in (a,b)$,
dan \lemmaref{relminmax:lemma} berlaku: Karena $g$ mencapai maksimum
di $c$, kita peroleh $g'(c) = 0$
dan karenanya $f'(c) = y$.

Demikian pula, jika $f'(a) > y > f'(b)$, gunakan $g(x) \coloneqq f(x)- yx$.
\end{proof}

Kita telah melihat bahwa
terdapat fungsi diskontinu yang memiliki
sifat nilai antara.  Meskipun pada awalnya sulit dibayangkan,
terdapat pula fungsi yang diferensiabel di mana-mana tetapi turunannya tidak
kontinu.

\begin{example} \label{baddifffunc:example}
Misalkan $f \colon \R \to \R$ fungsi yang didefinisikan oleh
\begin{equation*}
f(x) \coloneqq
\begin{cases}
{\bigl( x \sin(\nicefrac{1}{x}) \bigr)}^2 & \text{jika } x \neq 0, \\
0 & \text{jika } x = 0.
\end{cases}
\end{equation*}
Kita klaim bahwa $f$ diferensiabel di mana-mana, tetapi
$f' \colon \R \to \R$ tidak kontinu di
titik asal.  Selain itu, $f$ memiliki minimum di $0$, tetapi turunannya
berganti tanda tak berhingga kali di dekat titik asal.
Lihat \figureref{fig:nonc1diff}.
\begin{myfigureht}
\myincludepdft{nonc1diff_full}{%
Di sebelah kiri terdapat grafik fungsi yang berosilasi antara
sumbu x dan fungsi x kuadrat, yang grafiknya digambarkan
sebagai garis putus-putus.
Di sebelah kanan terdapat fungsi yang berosilasi semakin cepat
ketika mendekati titik asal, agak menyerupai grafik
sinus dari 1 per x, tetapi seluruh grafik tampak sedikit dimiringkan
sehingga juga bergerak ke atas.}
\caption{Fungsi dengan turunan diskontinu. Fungsi $f$ berada di sebelah kiri
dan $f'$ di sebelah kanan.  Perhatikan bahwa $f(x) \leq x^2$ pada grafik kiri.\label{fig:nonc1diff}}
\end{myfigureht}

Bukti: Langsung dari definisi, $f$ memiliki
minimum mutlak di $0$; kita mengetahui $f(x) \geq 0$ untuk semua $x$ dan $f(0) = 0$.

Untuk $x \neq 0$,
$f$ diferensiabel
dan
turunannya
adalah $2 \sin (\nicefrac{1}{x}) \bigl( x \sin (\nicefrac{1}{x}) -
\cos(\nicefrac{1}{x}) \bigr)$.
Sebagai latihan, tunjukkan bahwa untuk $x_n = \frac{4}{(8n+1)\pi}$,
berlaku
$\lim_{n\to\infty} f'(x_n) = -1$, dan untuk
$y_n = \frac{4}{(8n+3)\pi}$, berlaku
$\lim_{n\to\infty} f'(y_n) = 1$.  Jadi, $f'$
tidak mungkin kontinu di $0$
apa pun nilai $f'(0)$.

Mari kita tunjukkan bahwa $f$ diferensiabel di $0$ dan $f'(0)=0$.
Untuk $x \neq 0$,
\begin{equation*}
\abs{\frac{f(x)-f(0)}{x-0} - 0}
=
\abs{\frac{x^2 \sin^2(\nicefrac{1}{x})}{x}}
=
\babs{x \sin^2(\nicefrac{1}{x})}
\leq
\sabs{x} .
\end{equation*}
Tentu saja, ketika $x$ menuju nol, $\sabs{x}$ menuju nol, sehingga
$\abs{\frac{f(x)-f(0)}{x-0} - 0}$ menuju nol.  Oleh karena itu, $f$
diferensiabel di 0 dan turunannya di 0 bernilai 0.
Salah satu hal penting dalam perhitungan di atas
adalah $\babs{f(x)} \leq x^2$;
lihat juga Latihan \ref{exercise:bndmuldiff} dan
\ref{exercise:diffsqueeze}.
\end{example}

Kadang-kadang berguna untuk mengandaikan bahwa turunan dari suatu fungsi
diferensiabel bersifat kontinu.  Jika $f \colon I \to \R$ diferensiabel dan
turunannya $f'$ kontinu pada $I$, kita katakan bahwa $f$
\emph{\myindex{diferensiabel secara
kontinu}}\index{diferensiabel!secara kontinu}.  Lazim ditulis
$C^1(I)$ untuk himpunan fungsi yang diferensiabel secara kontinu
pada $I$.

\subsection{Latihan}

\begin{exercise}
Selesaikan pembuktian \propref{incdecdiffprop}.
\end{exercise}

\begin{exercise}
Selesaikan pembuktian \propref{firstderminmaxtest}.
\end{exercise}

\begin{exercise} \label{exercise:boundeddermeanslip}
Misalkan $f \colon \R \to \R$ adalah fungsi diferensiabel
sedemikian sehingga $f'$ merupakan fungsi terbatas.  Buktikan
bahwa $f$ merupakan fungsi kontinu Lipschitz.
\end{exercise}

\begin{exercise}
\pagebreak[1]
Misalkan $f \colon [a,b] \to \R$ diferensiabel dan $c \in [a,b]$.
Tunjukkan bahwa terdapat barisan $\{ x_n \}_{n=1}^\infty$ yang konvergen menuju $c$, dengan $x_n
\neq c$ untuk setiap $n$, sedemikian sehingga
\begin{equation*}
f'(c) = \lim_{n\to \infty} f'(x_n).
\avoidbreak
\end{equation*}
Perhatikan bahwa hal ini \emph{tidak} menyiratkan bahwa $f'$ kontinu (mengapa?).
\end{exercise}

\begin{exercise}
Misalkan $f \colon \R \to \R$ adalah fungsi sedemikian sehingga
$\babs{f(x)-f(y)} \leq \sabs{x-y}^2$ untuk setiap $x$ dan $y$.  Tunjukkan bahwa
$f(x) = C$ untuk suatu konstanta $C$.  Petunjuk: Tunjukkan bahwa $f$ diferensiabel
di setiap titik dan hitunglah turunannya.
\end{exercise}

\begin{exercise} \label{exercise:posderincr}
Selesaikan pembuktian \propref{incdecdiffstrictprop}.  Artinya,
misalkan $I$ adalah interval dan
$f \colon I \to \R$ adalah fungsi diferensiabel sedemikian sehingga
$f'(x) > 0$ untuk setiap $x \in I$.
Tunjukkan bahwa $f$ naik tegas.
\end{exercise}

\begin{exercise}
Misalkan $f \colon (a,b) \to \R$ adalah fungsi diferensiabel
sedemikian sehingga
$f'(x) \neq 0$ untuk setiap $x \in (a,b)$.  Misalkan
terdapat
titik $c \in (a,b)$ sedemikian sehingga $f'(c) > 0$.
Buktikan bahwa $f'(x) > 0$ untuk setiap $x \in (a,b)$.
\end{exercise}

\begin{exercise} \label{exercise:samediffconst}
Misalkan $f \colon (a,b) \to \R$ dan $g \colon (a,b) \to \R$ adalah
fungsi-fungsi diferensiabel sedemikian sehingga $f'(x) = g'(x)$ untuk setiap $x \in (a,b)$;
tunjukkan bahwa terdapat konstanta $C$ sedemikian sehingga $f(x) = g(x) + C$.
\end{exercise}

\begin{exercise}
Buktikan versi \myindex{aturan L'H\^opital}\index{aturan L'Hospital} berikut.
Misalkan 
$f \colon (a,b) \to \R$ dan $g \colon (a,b) \to \R$ adalah fungsi-fungsi
diferensiabel dan $c \in (a,b)$.  Misalkan $f(c) = 0$, $g(c)=0$,
$g'(x) \neq 0$ apabila $x \neq c$, dan
limit $\nicefrac{f'(x)}{g'(x)}$ saat $x$ menuju $c$ ada.  Tunjukkan bahwa
\begin{equation*}
\lim_{x \to c} \frac{f(x)}{g(x)} = 
\lim_{x \to c} \frac{f'(x)}{g'(x)} .
\end{equation*}
Bandingkan dengan \exerciseref{exercise:simpleLHopital}.
Catatan: Sebelum melakukan hal lain, buktikan bahwa $g(x) \neq 0$ apabila $x \neq c$.
\end{exercise}

\begin{exercise}
Misalkan $f \colon (a,b) \to \R$ adalah fungsi diferensiabel tak terbatas.  Tunjukkan bahwa
$f' \colon (a,b) \to \R$ tak terbatas.
Catatan: Penting bahwa $(a,b)$ merupakan interval terbatas.
\end{exercise}

\begin{exercise}
Buktikan teorema yang sebenarnya dibuktikan Rolle pada tahun 1691:
\emph{Jika $f$ adalah polinomial,
$f'(a) = f'(b) = 0$ untuk suatu $a < b$,
dan tidak terdapat $c \in (a,b)$ sedemikian sehingga $f'(c) = 0$,
maka terdapat paling banyak satu akar $f$ dalam $(a,b)$,
yaitu paling banyak satu $x \in (a,b)$ sedemikian sehingga $f(x) = 0$.}
Dengan kata lain, di antara setiap dua akar berurutan dari $f'$ terdapat paling banyak satu
akar $f$.
Petunjuk: Misalkan terdapat dua akar dan lihatlah apa yang terjadi.
\end{exercise}

\begin{exercise}
Misalkan $a,b \in \R$ dan $f \colon \R \to \R$ diferensiabel,
$f'(x) = a$ untuk setiap $x$, dan $f(0) = b$.  Carilah $f$ dan buktikan bahwa 
fungsi tersebut adalah satu-satunya fungsi diferensiabel dengan sifat ini.
\end{exercise}

\begin{exercise} \label{exercise:endpointderivative}
\leavevmode
\begin{enumerate}[a)]
\item
Buktikan \propref{prop:endpointderivative}.
\item 
Misalkan $f \colon (a,b) \to \R$ kontinu, dan misalkan $f$ diferensiabel di setiap titik kecuali di $c \in
(a,b)$ serta $\lim_{x \to c} f'(x) = L$.
Buktikan bahwa $f$ diferensiabel di
$c$ dan $f'(c) = L$.
\end{enumerate}
\end{exercise}

\begin{exercise} \label{exercise:extendboundedder}
Misalkan $f \colon (0,1) \to \R$ diferensiabel dan $f'$
terbatas.
\begin{enumerate}[a)]
\item
Tunjukkan bahwa terdapat fungsi kontinu $g \colon [0,1) \to \R$
sedemikian sehingga $f(x) = g(x)$ untuk setiap $x \neq 0$.\\
Petunjuk: \propref{context:prop} dan
\exerciseref{exercise:boundeddermeanslip}.
\item
Carilah contoh dengan $g$ tidak diferensiabel di $x=0$.
\\
Petunjuk: Pertimbangkan sesuatu yang didasarkan pada $\sin(\ln x)$,
dan anggaplah bahwa sifat-sifat dasar
$\sin$ dan $\ln$ dari kalkulus telah diketahui.
\item
Alih-alih mengasumsikan bahwa $f'$ terbatas, asumsikan bahwa $\lim_{x \to 0} f'(x)
= L$.  Buktikan bahwa $g$ tidak hanya ada, tetapi juga diferensiabel di $0$
dan $g'(0) = L$.
\end{enumerate}
\end{exercise}

\begin{exercise}
Buktikan \thmref{thm:cauchymvt}.
\end{exercise}

%%%%%%%%%%%%%%%%%%%%%%%%%%%%%%%%%%%%%%%%%%%%%%%%%%%%%%%%%%%%%%%%%%%%%%%%%%%%%%

\sectionnewpage
\section{Teorema Taylor}
\label{sec:taylor}

%mbxINTROSUBSECTION

\sectionnotes{kurang dari satu kuliah (bagian opsional)}

\subsection{Turunan berorde lebih tinggi}

Jika $f \colon I \to \R$ diferensiabel, kita memperoleh fungsi
$f' \colon I \to \R$.  Fungsi
$f'$ disebut \emph{\myindex{turunan pertama}} dari $f$.
Jika $f'$ diferensiabel, dengan
$f'' \colon I \to \R$ kita menyatakan turunan dari $f'$.  Fungsi $f''$
disebut \emph{\myindex{turunan kedua}} dari $f$.
\glsadd{not:secondthirdfourthder}
Dengan cara yang sama, kita memperoleh
$f'''$, $f''''$, dan seterusnya.
Untuk turunan yang lebih banyak,
notasi tersebut menjadi tidak praktis; dengan
$f^{(n)}$ kita menyatakan
\glsadd{not:nthder}
\emph{turunan ke-$n$}\index{turunan ke-n@turunan ke-$n$} dari $f$.
%
Jika $f$ memiliki turunan hingga orde ke-$n$, kita mengatakan bahwa $f$
\emph{diferensiabel $n$ kali}\index{diferensiabel n kali@diferensiabel $n$ kali}
\index{diferensiabel!n kali@$n$ kali}.

\subsection{Teorema Taylor}

Teorema Taylor%
\footnote{Dinamai menurut matematikawan Inggris
\href{https://en.wikipedia.org/wiki/Brook_Taylor}{Brook Taylor}
(1685--1731).
Teorema ini ditemukan oleh
matematikawan Skotlandia
\href{https://en.wikipedia.org/wiki/James_Gregory_(mathematician)}{James Gregory}
(1638--1675).  Pernyataan yang kami berikan
dikemukakan oleh
\href{https://en.wikipedia.org/wiki/Lagrange}{Joseph-Louis Lagrange}
(1736--1813).} (setidaknya versi yang kami berikan di sini)
merupakan generalisasi dari \hyperref[thm:mvt]{teorema nilai rata-rata}.
Teorema nilai rata-rata menyatakan bahwa, dengan suatu galat kecil, $f(x)$ untuk $x$ di dekat $x_0$ dapat
diaproksimasi oleh $f(x_0)$:
\begin{equation*}
f(x) = f(x_0) + f'(c)(x-x_0),
\end{equation*}
di mana \myquote{galat} $f'(c)(x-x_0)$
dinyatakan dalam turunan pertama
di suatu titik $c$ antara $x$ dan $x_0$.
Teorema Taylor menggeneralisasi hasil ini ke turunan berorde lebih tinggi.
Teorema ini menyatakan bahwa, dengan suatu galat kecil, setiap fungsi yang
diferensiabel $n$ kali dapat diaproksimasi di suatu titik $x_0$
oleh polinom berderajat $n$.  Galat
aproksimasi ini menuju nol lebih cepat daripada ${(x-x_0)}^{n}$ ketika
$x$ menuju $x_0$.
Untuk melihat mengapa ini merupakan aproksimasi yang baik, perhatikan bahwa untuk $n$ yang besar, 
${(x-x_0)}^n$ sangat kecil dalam interval kecil di sekitar $x_0$.
Kita akan memerlukan satu turunan lagi (jadi $n+1$) untuk menuliskan galat
seperti yang kita lakukan dalam teorema nilai rata-rata.

\begin{defn}
Untuk suatu fungsi yang diferensiabel $n$ kali, katakanlah $f$, yang didefinisikan di sekitar titik $x_0 \in \R$, kita mendefinisikan
\emph{\myindex{polinom Taylor}} berorde ke-$n$%
\index{polinom Taylor berorde ke-n@polinom Taylor berorde ke-$n$}
untuk $f$ di $x_0$ sebagai
\begin{equation*}
\begin{split}
P_n^{x_0}(x)
& \coloneqq
\sum_{k=0}^n
\frac{f^{(k)}(x_0)}{k!}{(x-x_0)}^k
\\
& =
f(x_0)
+ f'(x_0)(x-x_0)
+ \frac{f''(x_0)}{2}{(x-x_0)}^2
+ \frac{f^{(3)}(x_0)}{6}{(x-x_0)}^3
\\
& \qquad
+ \cdots
+ \frac{f^{(n)}(x_0)}{n!}{(x-x_0)}^n .
\end{split}
\end{equation*}
\end{defn}

Lihat \figureref{fig:taylorsin} untuk
polinom Taylor berderajat ganjil bagi fungsi sinus di $x_0=0$.
Semua suku berderajat genap bernilai nol, sebab turunan berorde genap 
dari sinus kembali berupa sinus, yang bernilai nol di titik asal.
\begin{myfigureht}
\myincludegraphics{taylorsin}{%
Diberikan grafik fungsi sinus.  Ditampilkan pula beberapa
aproksimasi.  Pertama, sebuah garis lurus yang
menyinggung di titik asal diberi label y sama dengan P sub 1 super 0 dari x.
Kedua, grafik sebuah polinom kubik yang menyinggung grafik sinus
di titik asal dan cukup mengaproksimasinya di sekitar titik tersebut diberi label
y sama dengan P sub 3 super 0 dari x.
Demikian pula untuk aproksimasi berderajat 5 dan berderajat 7.  Semakin tinggi
derajatnya, semakin jauh dari titik asal grafik tersebut masih cukup dekat 
dengan fungsi sinus; sebagai contoh, aproksimasi berderajat 7 tampak
cukup dekat di antara dua titik tempat sinus
bernilai nol di sekitar titik asal, yaitu minus dan plus pi.}
\caption{Polinom Taylor berderajat ganjil untuk fungsi
sinus.\label{fig:taylorsin}}
\end{myfigureht}

Pernyataan teorema Taylor yang kami berikan mencakup
\hyperref[thm:mvt]{teorema nilai rata-rata}, yang dapat dipandang sebagai
teorema Taylor untuk $n=0$.

\begin{thm}[Taylor]\index{teorema Taylor} \label{thm:taylor}
Misalkan $f \colon [a,b] \to \R$ adalah fungsi dengan $n$ turunan yang
kontinu pada $[a,b]$ dan sedemikian sehingga $f^{(n+1)}$ ada pada $(a,b)$.
Untuk titik-titik berbeda $x_0$ dan $x$ di $[a,b]$,
terdapat titik $c$ antara $x_0$
dan $x$ sedemikian sehingga
\begin{equation*}
f(x)=P_{n}^{x_0}(x)+\frac{f^{(n+1)}(c)}{(n+1)!}{(x-x_0)}^{n+1} .
\end{equation*}
\end{thm}

Suku $R_n^{x_0}(x)\coloneqq
f(x)-P_n^{x_0}(x) =
\frac{f^{(n+1)}(c)}{(n+1)!}{(x-x_0)}^{n+1}$ disebut
\emph{suku sisa}\index{suku sisa dalam rumus Taylor}.  Bentuk ini
disebut
\emph{\myindex{bentuk Lagrange}} dari suku sisa.  Ada cara lain
untuk menuliskan suku sisa, tetapi tidak kita bahas.  Perhatikan bahwa $c$ bergantung pada
$x$ dan $x_0$.

\begin{proof}
Misalkan $M_{x,x_0}$ adalah bilangan (yang bergantung pada $x$ dan $x_0$) yang memenuhi persamaan
\begin{equation*}
f(x)=P_{n}^{x_0}(x)+M_{x,x_0}{(x-x_0)}^{n+1} .
\end{equation*}
Definisikan fungsi $g(s)$ dengan
\begin{equation*}
g(s) \coloneqq f(s)-P_n^{x_0}(s)-M_{x,x_0}{(s-x_0)}^{n+1} .
\end{equation*}
Kita menghitung
turunan ke-$k$ polinom Taylor di $x_0$ dan memperoleh
${(P_n^{x_0})}^{(k)}(x_0) = f^{(k)}(x_0)$ untuk
$k=0,1,2,\ldots,n$ (turunan ke-nol dari suatu fungsi adalah fungsi itu
sendiri).  Oleh karena itu,
\begin{equation*}
g(x_0) = g'(x_0) = g''(x_0) = \cdots = g^{(n)}(x_0) = 0 .
\end{equation*}
Khususnya, $g(x_0) = 0$.
Kita juga memiliki $g(x) = 0$.
Berdasarkan
\hyperref[thm:mvt]{teorema nilai rata-rata},
terdapat $x_1$ antara $x_0$ dan $x$ sedemikian sehingga $g'(x_1) = 0$.
Dengan menerapkan \hyperref[thm:mvt]{teorema nilai rata-rata}
pada $g'$, kita memperoleh bahwa terdapat
$x_2$ antara $x_0$ dan $x_1$ (dan karena itu antara $x_0$ dan $x$)
sedemikian sehingga $g''(x_2) = 0$.  Kita mengulangi
argumen ini $n+1$ kali untuk memperoleh bilangan $x_{n+1}$ antara $x_0$ dan $x_n$
(dan karena itu antara $x_0$ dan $x$) sedemikian sehingga $g^{(n+1)}(x_{n+1}) = 0$.

Misalkan $c \coloneqq x_{n+1}$.
Kita menghitung turunan ke-$(n+1)$ dari $g$ dan memperoleh
\begin{equation*}
g^{(n+1)}(s) = f^{(n+1)}(s)-(n+1)!\,M_{x,x_0} .
\end{equation*}
Dengan menyubstitusikan $c$ untuk $s$, kita memperoleh $M_{x,x_0} = \frac{f^{(n+1)}(c)}{(n+1)!}$, dan
pembuktian selesai.
\end{proof}

Dalam pembuktian tersebut, kita memperoleh
${(P_n^{x_0})}^{(k)}(x_0) = f^{(k)}(x_0)$ untuk $k=0,1,2,\ldots,n$.
Oleh karena itu, polinom Taylor memiliki turunan yang sama dengan $f$ di $x_0$
hingga turunan ke-$n$.  Itulah sebabnya polinom Taylor merupakan
aproksimasi yang baik bagi $f$.
Perhatikan dalam \figureref{fig:taylorsin} bahwa polinom-polinom Taylor merupakan
aproksimasi yang cukup baik bagi sinus di dekat $x=0$.

Polinom Taylor tidak selalu memberikan aproksimasi yang baik
di setiap titik.
Pertimbangkan pengembangan fungsi
$f(x) \coloneqq \frac{x}{1-x}$ di sekitar 0.
Untuk $x < 1$, kita memperoleh grafik-grafik dalam
\figureref{fig:taylorgeom}.  Garis-garis titik-titik adalah aproksimasi berderajat
pertama, kedua, dan ketiga.
Garis putus-putus adalah polinom berderajat 20.
Aproksimasi tersebut memang tampak menjadi
lebih baik seiring naiknya derajat untuk $x > -1$.
Untuk $x < -1$, aproksimasi itu justru tampak semakin buruk.
Polinom-polinom tersebut
adalah jumlah parsial deret geometri $\sum_{n=1}^\infty x^n$,
dan deret tersebut hanya konvergen pada $(-1,1)$.
Lihat pembahasan deret pangkat dalam
\sectionref{sec:moreonseries}.

\begin{myfigureht}
\myincludegraphics{taylorgeom}{%
Dengan garis tebal ditampilkan grafik suatu fungsi pada interval
dari minus 2 hingga 1, yang memasuki gambar secara mendatar dari kiri di bawah
sumbu-x, memotong sumbu-x di titik asal, lalu berbelok tajam
ke atas dan meninggalkan gambar sebelum mencapai 1.
Grafik tiga polinom Taylor pertama ditampilkan dengan
garis titik-titik.  Polinom-polinom itu berupa
garis lurus, polinom kuadrat, dan polinom kubik.
Ketiganya tampak menjadi 
aproksimasi yang cukup wajar tepat di sekitar titik asal,
tetapi memburuk ketika kita menjauh.
Digambarkan sebuah garis putus-putus yang menandakan aproksimasi berderajat 20, dan
dari suatu titik yang sedikit lebih besar daripada minus 1 hingga tempat garis
tebal meninggalkan gambar di sebelah kanan, perbedaan antara keduanya
tidak terlihat.  Namun, tepat di sekitar minus 1,
jika kita bergerak dari kanan ke kiri, garis itu melesat naik meninggalkan gambar,
sehingga untuk x yang lebih kecil daripada minus 1 aproksimasi tersebut jauh lebih buruk
daripada aproksimasi berderajat rendah (yang sejak awal sudah tidak baik).}
\caption{Fungsi $\frac{x}{1-x}$, polinom Taylor
$P_1^0$, $P_2^0$, $P_3^0$ (semuanya bergaris titik-titik), dan polinom $P_{20}^0$
(bergaris putus-putus).\label{fig:taylorgeom}}
\end{myfigureht}

Jika $f$ \emph{\myindex{diferensiabel
tak berhingga kali}}\index{diferensiabel!tak berhingga kali},
yaitu jika $f$ dapat
didiferensialkan sebanyak berapa pun, maka 
kita mendefinisikan \emph{\myindex{deret Taylor}}:
\begin{equation*}
%T^{x_0}(x)
%\coloneqq
\sum_{k=0}^\infty
\frac{f^{(k)}(x_0)}{k!}{(x-x_0)}^k .
\end{equation*}
Tidak ada jaminan bahwa deret ini konvergen untuk sembarang
$x \neq x_0$.  Bahkan di tempat deret itu konvergen, tidak ada jaminan
bahwa deret itu konvergen ke fungsi $f$.  Fungsi $f$
yang deret Taylornya di setiap titik $x_0$
konvergen ke $f$ pada suatu interval terbuka yang memuat $x_0$
disebut
\emph{fungsi analitik}\index{fungsi analitik}.
Banyak fungsi yang lazim dijumpai dalam praktik bersifat analitik.
Lihat \exerciseref{exercise:nonanalytic} untuk contoh fungsi
nonanalitik.

\medskip

Definisi turunan menyatakan bahwa
suatu fungsi
diferensiabel jika fungsi tersebut
secara lokal diaproksimasi oleh sebuah garis.
Sebagai catatan, terdapat suatu kebalikan dari teorema
Taylor,
yang tidak akan kita nyatakan ataupun buktikan,
yang mengatakan bahwa jika suatu fungsi
secara lokal diaproksimasi dengan cara tertentu oleh polinom berderajat $d$, maka fungsi itu
memiliki $d$ turunan.

\medskip

Teorema Taylor memberikan pembuktian singkat untuk suatu versi
uji turunan kedua.
Yang dimaksud dengan \emph{\myindex{minimum relatif ketat}}\index{minimum!relatif ketat}
dari $f$ di $c$
adalah bahwa terdapat $\delta > 0$ sedemikian sehingga $f(x) > f(c)$ untuk
semua $x \in (c-\delta,c+\delta)$ dengan $x\neq c$.
\emph{\myindex{Maksimum relatif ketat}}\index{maksimum!relatif ketat}
didefinisikan dengan cara serupa.
Kekontinuan turunan kedua tidak diperlukan, tetapi pembuktiannya lebih
sulit dan diserahkan sebagai latihan.  Pembuktian itu juga dapat digeneralisasi
menjadi uji turunan ke-$n$, yang juga diserahkan sebagai
latihan.

\begin{prop}[Uji turunan kedua]\index{uji turunan kedua}
Misalkan $f \colon (a,b) \to \R$ diferensiabel dua kali secara kontinu,
$x_0 \in (a,b)$, $f'(x_0) = 0$, dan $f''(x_0) > 0$.  Maka $f$ memiliki minimum relatif
ketat di $x_0$.
\end{prop}

\begin{proof}
Karena $f''$ kontinu, terdapat $\delta > 0$
sedemikian sehingga $f''(c) > 0$ untuk semua $c \in (x_0-\delta,x_0+\delta)$;
lihat \exerciseref{exercise:positivecontneigh}.
Ambil $x \in (x_0-\delta,x_0+\delta)$, $x \neq x_0$.
Teorema Taylor menyatakan bahwa untuk suatu $c$ antara $x_0$ dan $x$,
\begin{equation*}
f(x) 
=
f(x_0) + f'(x_0) (x-x_0) +
\frac{f''(c)}{2}{(x-x_0)}^{2} 
=
f(x_0) + \frac{f''(c)}{2}{(x-x_0)}^{2}  .
\end{equation*}
Karena $f''(c) > 0$ dan ${(x-x_0)}^2 > 0$, kita memperoleh $f(x) > f(x_0)$.
\end{proof}

\subsection{Latihan}

\begin{exercise}
Hitunglah polinom Taylor ke-$n$ di $0$ untuk fungsi eksponensial.
\end{exercise}

\begin{exercise}
Misalkan $p$ adalah polinomial berderajat $d$.  Untuk $x_0 \in \R$,
tunjukkan bahwa
polinom Taylor ke-$d$ bagi $p$ di $x_0$ sama dengan $p$.
\end{exercise}

\begin{exercise}
Misalkan $f(x) \coloneqq \sabs{x}^3$.  Hitunglah $f'(x)$ dan $f''(x)$ untuk setiap $x$,
tetapi tunjukkan bahwa $f^{(3)}(0)$ tidak ada.
\end{exercise}

\begin{exercise}
Misalkan $f \colon \R \to \R$ memiliki $n$ turunan kontinu.  Tunjukkan
bahwa untuk setiap $x_0 \in \R$,
terdapat polinomial $P$ dan $Q$ berderajat $n$ serta 
suatu $\epsilon > 0$ sedemikian sehingga $P(x) \leq f(x) \leq Q(x)$ untuk setiap $x \in
[x_0,x_0+\epsilon]$  dan
$Q(x)-P(x) = \lambda {(x-x_0)}^n$ untuk suatu $\lambda \geq 0$.
\end{exercise}

\begin{exercise}
Jika $f \colon [a,b] \to \R$ memiliki $n+1$ turunan kontinu\footnote{%
Pernyataan yang sama berlaku jika kita hanya mensyaratkan $n$ turunan,
tetapi pembuktiannya jauh lebih sulit.
}
dan $x_0 \in [a,b]$,
buktikan
$\lim\limits_{x\to x_0} \frac{R_n^{x_0}(x)}{{(x-x_0)}^n} = 0$.
\end{exercise}

\begin{exercise}
Misalkan $f \colon [a,b] \to \R$ memiliki $n+1$ turunan kontinu
dan $x_0 \in (a,b)$.
Buktikan: $f^{(k)}(x_0) = 0$ untuk setiap $k = 0, 1, 2, \ldots, n$
jika dan hanya jika $\lim\limits_{x\to x_0} \frac{f(x)}{{(x-x_0)}^{n+1}}$ ada.
\end{exercise}

\begin{exercise}
Misalkan $a,b,c \in \R$ dan $f \colon \R \to \R$ diferensiabel dua kali,
$f''(x) = a$ untuk setiap $x$, $f'(0) = b$, dan $f(0) = c$.  Carilah $f$ dan buktikan bahwa 
fungsi tersebut adalah satu-satunya fungsi diferensiabel dengan sifat-sifat ini.
\end{exercise}

\begin{exercise}[Menantang]
Tunjukkan bahwa suatu konvers sederhana dari teorema Taylor tidak berlaku.
Carilah fungsi $f \colon \R \to \R$ yang tidak memiliki turunan kedua di $x=0$ sedemikian sehingga
$\babs{f(x)} \leq \babs{x^3}$, yakni, $f$ mendekati nol di 0 lebih cepat daripada $x^2$, dan
meskipun $f'(0)$ ada, $f''(0)$ tidak ada.
\end{exercise}

\begin{exercise} \label{exercise:extendboundedder2}
Misalkan $f \colon (0,1) \to \R$ diferensiabel dan $f''$
terbatas.
\begin{enumerate}[a)]
\item
Tunjukkan bahwa terdapat fungsi $g \colon [0,1) \to \R$ yang diferensiabel satu kali
sedemikian sehingga $f(x) = g(x)$ untuk setiap $x \neq 0$.  Petunjuk: 
Lihat
\exerciseref{exercise:extendboundedder}.
\item
Carilah contoh dengan $g$ yang tidak diferensiabel dua kali di $x=0$.
\end{enumerate}
\end{exercise}

\begin{exercise}
Buktikan
\emph{uji turunan ke-$n$}\index{uji turunan ke-n@uji turunan ke-$n$}.
Misalkan $n \in \N$,
$x_0 \in (a,b)$, dan $f \colon (a,b) \to \R$ diferensiabel secara kontinu hingga orde ke-$n$,
dengan $f^{(k)}(x_0) = 0$ untuk $k=1,2,\ldots,n-1$, dan
$f^{(n)}(x_0)
\neq 0$.
Buktikan:
\begin{enumerate}[a)]
\item
Jika $n$ ganjil, maka $f$ tidak memiliki minimum relatif
maupun maksimum relatif di $x_0$.
\item
Jika $n$ genap, maka $f$ memiliki minimum relatif tegas di $x_0$ jika
$f^{(n)}(x_0) > 0$
dan maksimum relatif tegas di $x_0$ jika $f^{(n)}(x_0) < 0$.
\end{enumerate}
\end{exercise}

\begin{exercise}
Buktikan versi yang lebih umum dari uji turunan kedua.
Misalkan $f \colon (a,b) \to \R$ diferensiabel dan $x_0 \in (a,b)$
sedemikian sehingga $f'(x_0) = 0$, $f''(x_0)$ ada, dan $f''(x_0) > 0$.
Buktikan bahwa $f$ memiliki minimum
relatif tegas di $x_0$.  Petunjuk: Tinjaulah definisi limit bagi $f''(x_0)$.
\end{exercise}

%%%%%%%%%%%%%%%%%%%%%%%%%%%%%%%%%%%%%%%%%%%%%%%%%%%%%%%%%%%%%%%%%%%%%%%%%%%%%%

\sectionnewpage
\section{Teorema fungsi invers}
\label{sec:ift}

%mbxINTROSUBSECTION

\sectionnotes{kurang dari 1 pertemuan (bagian opsional, diperlukan untuk
\sectionref{sec:logandexp}, memerlukan
\sectionref{sec:monotonefunc})}

\subsection{Teorema fungsi invers}

Kita mulai dengan contoh sederhana.  Perhatikan fungsi $f(x) \coloneqq a x$ untuk suatu
bilangan $a \neq 0$.  Fungsi $f \colon \R \to \R$ bersifat bijektif, dan inversnya
adalah $f^{-1}(y) = \frac{1}{a} y$.  Secara khusus, $f'(x) = a$ dan
$(f^{-1})'(y) = \frac{1}{a}$.  Karena fungsi diferensiabel
\myquote{secara infinitesimal menyerupai} fungsi linear, kita mengharapkan
perilaku serupa dari invers suatu fungsi diferensiabel.
Gagasan utama untuk mendiferensialkan fungsi invers diberikan oleh lema berikut.

\begin{lemma} \label{lemma:ift}
Misalkan $I,J \subset \R$ merupakan interval.
Jika $f \colon I \to J$ monoton tegas (dan karenanya injektif),
surjektif ($f(I) = J$),
diferensiabel di $x_0 \in I$, dan $f'(x_0) \neq 0$,
maka invers
$f^{-1}$ diferensiabel di $y_0 = f(x_0)$ dan
\begin{equation*}
(f^{-1})'(y_0) = \frac{1}{f'\bigl( f^{-1}(y_0) \bigr)} = \frac{1}{f'(x_0)} .
\end{equation*}
Jika $f$ diferensiabel secara kontinu dan $f'$ tidak pernah nol, maka $f^{-1}$
diferensiabel secara kontinu.
\end{lemma}

\begin{proof}
Berdasarkan \propref{prop:invcont}, $f$ memiliki invers yang kontinu.  Untuk memudahkan,
namakan invers tersebut $g \colon J \to I$.
Misalkan $x_0,y_0$ seperti pada pernyataan.  Untuk $x \in I$, tuliskan $y \coloneqq f(x)$.
Jika $x \neq x_0$, dan dengan demikian $y \neq y_0$, kita peroleh
\begin{equation*}
\frac{g(y)-g(y_0)}{y-y_0} =
\frac{g\bigl(f(x)\bigr)-g\bigl(f(x_0)\bigr)}{f(x)-f(x_0)} =
\frac{x-x_0}{f(x)-f(x_0)} .
\end{equation*}
Lihat \figureref{inversefig} untuk gagasan geometrisnya.
\begin{myfigureht}
\myincludepdft{inversefigAB}{%
Dua diagram.
Di sebelah kiri, suatu fungsi f yang menyerupai fungsi akar kuadrat
ditampilkan pada kuadran pertama.  Pada sumbu horizontal
terdapat titik berlabel x sama dengan g dari y.  Sebuah garis putus-putus
bergerak vertikal ke atas dari titik ini hingga mencapai titik yang bersesuaian
pada grafik f.  Garis putus-putus lain, kini horizontal, bergerak ke kiri
hingga mencapai sumbu y dan titik yang ditandai f dari x sama dengan y.
Titik lain, x subskrip nol sama dengan g dari y subskrip nol, juga ditandai
sedikit di sebelah kanan titik pertama.  Susunan yang sama dengan
garis-garis putus-putus diulang, tetapi x diganti dengan x subskrip nol dan
y diganti dengan y subskrip nol.
Sebuah garis sekan digambar melalui kedua titik pada grafik
dan diberi label dengan kemiringan
f dari x dikurangi f dari x subskrip nol, seluruhnya dibagi besaran
x dikurangi x subskrip nol,
yang sama dengan
y dikurangi y subskrip nol, seluruhnya dibagi besaran
g dari y dikurangi g dari y subskrip nol.
Diagram di sebelah kanan sama, tetapi gambar
dicerminkan terhadap garis x sama dengan y, sehingga sumbu x
menjadi sumbu y dan sebaliknya, dan grafik kini lebih menyerupai
x kuadrat.  Kedua bentuk kemiringan garis sekan tersebut
saling berkebalikan.}
\caption{Interpretasi turunan fungsi
invers.\label{inversefig}}
\end{myfigureht}

Misalkan
\begin{equation*}
Q(x) \coloneqq
\begin{cases}
\frac{x-x_0}{f(x)-f(x_0)} & \text{jika } x \neq x_0, \\
\frac{1}{f'(x_0)}         & \text{jika } x = x_0 \quad \text{(perhatikan bahwa }
f'(x_0) \neq 0 \text{)}.
\end{cases}
\end{equation*}
Karena $f$ diferensiabel di $x_0$,
\begin{equation*}
\lim_{x \to x_0} Q(x) =
\lim_{x \to x_0} 
\frac{x-x_0}{f(x)-f(x_0)} 
=
\frac{1}{f'(x_0)}
=
Q(x_0) ,
\end{equation*}
dengan kata lain, $Q$ kontinu di $x_0$.
Karena $g$ kontinu di $y_0$,
komposisi $Q\bigl(g(y)\bigr)$ kontinu di $y_0$ berdasarkan
\propref{prop:compositioncont}.  Untuk $y \neq y_0$, berlaku
$Q\bigl(g(y)\bigr) = \frac{g(y)-g(y_0)}{y-y_0}$.
Oleh karena itu,
\begin{equation*}
\frac{1}{f'\bigl(g(y_0)\bigr)}
= Q\bigl(g(y_0)\bigr)
= \lim_{y \to y_0} Q\bigl(g(y)\bigr)
= \lim_{y \to y_0} \frac{g(y)-g(y_0)}{y-y_0} .
\end{equation*}
Jadi, $g$ diferensiabel di $y_0$, dan $g'(y_0) =
\frac{1}{f'\left(\vphantom{1^1_1}g(y_0)\right)}$.

Jika $f'$ kontinu dan tidak nol untuk semua $x \in I$,
maka lema tersebut berlaku untuk semua $x \in I$.  Karena $g$ juga
kontinu (fungsi ini diferensiabel), turunan $g'(y) =
\frac{1}{f'\left(\vphantom{1^1_1}g(y)\right)}$ pasti kontinu.
\end{proof}

Hasil berikut biasanya disebut teorema fungsi invers.

\begin{thm}[Teorema fungsi invers]\index{teorema fungsi invers}
Misalkan $f \colon (a,b) \to \R$ fungsi yang diferensiabel secara kontinu,
dan $x_0 \in (a,b)$ suatu titik dengan $f'(x_0) \neq 0$.  Maka terdapat
interval terbuka $I \subset (a,b)$ dengan $x_0 \in I$, sehingga
pembatasan $f|_{I}$ injektif dan memiliki invers yang diferensiabel secara kontinu,
$g \colon J \to I$, yang didefinisikan pada interval $J \coloneqq f(I)$,
serta
\begin{equation*}
g'(y) = \frac{1}{f'\bigl( g(y) \bigr)} \qquad \text{untuk semua } y \in J.
\end{equation*}
\end{thm}

\begin{proof}
Tanpa mengurangi keumuman, andaikan $f'(x_0) > 0$.  Karena $f'$
kontinu, pasti terdapat interval terbuka $I = (x_0-\delta,x_0+\delta)$
sedemikian sehingga $f'(x) > 0$ untuk semua $x \in I$.  Lihat
\exerciseref{exercise:positivecontneigh}.

Berdasarkan \propref{incdecdiffstrictprop}, $f$ naik tegas
pada $I$, sehingga pembatasan $f|_{I}$ bijektif ke $J \coloneqq f(I)$.
Karena $f$ kontinu,
\corref{cor:continterval}
(atau secara langsung melalui
\hyperref[IVT:thm]{teorema nilai antara})
menyiratkan bahwa
$f(I)$ merupakan interval.
Sekarang terapkan \lemmaref{lemma:ift}.
\end{proof}

Dalam \exampleref{example:sqrt2}, kita melihat betapa sulitnya
membuktikan keberadaan $\sqrt{2}$ tanpa alat apa pun.
Dengan \hyperref[IVT:thm]{teorema nilai antara},
keberadaan akar hampir sepele, dan
dengan perangkat pada bagian ini, kita akan membuktikan
lebih dari sekadar keberadaan.

\begin{cor}
Untuk $n \in \N$ dan $x \geq 0$, terdapat tepat satu
bilangan $y \geq 0$ (dinotasikan $x^{1/n} \coloneqq y$) sedemikian sehingga $y^n = x$.  Selain itu,
fungsi $g \colon (0,\infty) \to (0,\infty)$ yang didefinisikan oleh
$g(x) \coloneqq x^{1/n}$ diferensiabel secara kontinu dan
\begin{equation*}
g'(x) = \frac{1}{nx^{(n-1)/n}} = \frac{1}{n} \, x^{(1-n)/n} ,
\end{equation*}
dengan konvensi $x^{m/n} \coloneqq {(x^{1/n})}^{m}$.
\end{cor}

\begin{proof}
Untuk $x=0$, keberadaan tepat satu akar bersifat sepele.

Misalkan $f \colon (0,\infty) \to (0,\infty)$ didefinisikan oleh $f(y) \coloneqq y^n$.
Fungsi $f$ diferensiabel secara kontinu,
dan $f'(y) = ny^{n-1}$; lihat \exerciseref{exercise:diffofxn}.
Untuk $y > 0$, turunan $f'$ positif tegas,
sehingga, sekali lagi berdasarkan \propref{incdecdiffstrictprop}, $f$ naik
tegas (hal ini juga dapat dibuktikan secara langsung) dan karenanya injektif.
Andaikan $M$ dan $\epsilon$ sedemikian sehingga
$M > 1$ dan $1 > \epsilon > 0$.
Maka
$f(M) = M^n \geq M$ dan
$f(\epsilon) = \epsilon^n \leq \epsilon$.
Untuk setiap $x$ dengan $\epsilon < x < M$,
berdasarkan
\hyperref[IVT:thm]{teorema nilai antara}, kita peroleh $x \in
f\bigl( [\epsilon,M] \bigr) \subset
f\bigl( (0,\infty) \bigr)$.  Karena $M$ dan $\epsilon$ dipilih sembarang, $f$ surjektif ke
$(0,\infty)$, sehingga $f$ bijektif.
Misalkan $g$ invers dari $f$.  Kita memperoleh
keberadaan dan ketunggalan akar pangkat
$n$ yang positif.  \lemmaref{lemma:ift} mengatakan bahwa $g$ memiliki turunan
kontinu dan $g'(x) =
\frac{1}{f'\left(\vphantom{1^1_1}g(x)\right)} = \frac{1}{n {(x^{1/n})}^{n-1}}$.
\end{proof}

\begin{example}
Korolari tersebut memberikan contoh yang baik ketika teorema fungsi invers
memberi kita interval yang lebih kecil daripada $(a,b)$.  Ambil $f \colon \R \to \R$
yang didefinisikan oleh $f(x) \coloneqq x^2$.  Maka $f'(x_0) \neq 0$
selama $x_0 \neq 0$.  Jika $x_0 > 0$, kita dapat mengambil $I=(0,\infty)$, tetapi
tidak dapat mengambil interval yang lebih besar.
\end{example}

\begin{example}
Contoh lain yang berguna adalah $f(x) \coloneqq x^3$.  Fungsi $f \colon \R \to \R$
injektif dan surjektif, sehingga $f^{-1}(y) = y^{1/3}$ ada pada seluruh garis
bilangan real, termasuk nol dan $y$ negatif.  Fungsi $f$ memiliki
turunan kontinu, tetapi $f^{-1}$ tidak memiliki turunan di titik asal.  Alasannya
adalah $f'(0) = 0$.  Lihat grafiknya pada \figureref{cubecuberootfig}.
Perhatikan garis singgung vertikal pada grafik akar pangkat tiga di titik asal.
Lihat juga \exerciseref{exercise:oddroot}.
\begin{myfigureht}
\myincludegraphics{cubecuberoot}{%
Grafik y sama dengan akar pangkat tiga dari x digambar dengan garis tebal,
sedangkan grafik y sama dengan x pangkat tiga digambar dengan garis putus-putus.
Perhatikan bahwa meskipun grafik x pangkat tiga menyinggung sumbu x
di titik asal, grafik akar pangkat tiga merupakan pencerminannya terhadap
garis x sama dengan y dan oleh karena itu menyinggung sumbu y vertikal di titik asal.}
\caption{Grafik $x^3$ dan $x^{1/3}$.\label{cubecuberootfig}}
\end{myfigureht}
\end{example}


\subsection{Latihan}

\begin{exercise}
Andaikan $f \colon \R \to \R$ diferensiabel secara kontinu dan
$f'(x) > 0$ untuk semua $x$.  Tunjukkan bahwa $f$ memiliki invers pada interval $J =
f(\R)$, inversnya diferensiabel secara kontinu, dan ${(f^{-1})}'(y) >
0$ untuk semua $y \in f(\R)$.
\end{exercise}

\begin{exercise}
\pagebreak[2]
Andaikan $I,J$ merupakan interval dan fungsi monoton surjektif $f \colon I \to J$ memiliki invers $g \colon J \to I$.
Andaikan telah diketahui bahwa $f$ dan $g$ diferensiabel
di setiap titik dan $f'$ tidak pernah nol.  Dengan menggunakan aturan rantai tetapi bukan
\lemmaref{lemma:ift}, buktikan rumus
$g'(y) = \frac{1}{f'\left(\vphantom{1_1^1}g(y)\right)}$. % \bigl( \bigr) are overly big here!
Catatan: Latihan ini sama dengan \exerciseref{exercise:inversederformula};
Anda tidak perlu mengerjakannya lagi jika
sudah pernah menyelesaikannya.
\end{exercise}

\begin{exercise}
\pagebreak[2]
Misalkan $n\in \N$ genap.
Buktikan bahwa setiap $x > 0$ memiliki tepat satu akar pangkat $n$ yang negatif.
Dengan kata lain, terdapat bilangan negatif $y$ sedemikian sehingga $y^n = x$.
Hitung turunan
fungsi $g(x) \coloneqq y$.
\end{exercise}

\begin{exercise} \label{exercise:oddroot}
Misalkan $n \in \N$ ganjil dan $n \geq 3$.
Buktikan bahwa setiap $x$ memiliki tepat satu akar pangkat $n$.
Dengan kata lain, terdapat bilangan $y$ sedemikian sehingga $y^n = x$.  Buktikan bahwa
fungsi yang didefinisikan oleh $g(x) \coloneqq y$ diferensiabel kecuali di $x=0$
dan hitung turunannya.  Buktikan bahwa $g$ tidak diferensiabel di $x=0$.
\end{exercise}

\begin{exercise}[memerlukan \sectionref{sec:taylor}]
Tunjukkan bahwa jika dalam teorema fungsi invers $f$ memiliki $k$ turunan
kontinu, maka fungsi invers $g$ juga memiliki $k$ turunan
kontinu.
\end{exercise}

\begin{exercise}
Misalkan $f(x) \coloneqq x + 2 x^2 \sin(\nicefrac{1}{x})$ untuk $x \neq 0$ dan
$f(0) \coloneqq 0$.  Tunjukkan bahwa $f$ diferensiabel untuk setiap $x$, bahwa $f'(0) > 0$,
tetapi $f$ tidak memiliki invers
pada interval terbuka mana pun yang memuat titik asal.
\end{exercise}

\begin{exercise}
\leavevmode
\begin{enumerate}[a)]
\item
Misalkan $f \colon \R \to \R$ fungsi yang diferensiabel secara kontinu
dan $k > 0$ bilangan sedemikian sehingga $f'(x) \geq k$ untuk semua $x \in \R$.
Tunjukkan bahwa $f$ injektif dan surjektif, serta memiliki invers
$f^{-1} \colon \R \to \R$ yang diferensiabel secara kontinu.
\item
Temukan contoh $f \colon \R \to \R$
sedemikian sehingga $f'(x) > 0$
untuk semua $x$, tetapi $f$ tidak surjektif.
\end{enumerate}
\end{exercise}

\begin{exercise}
Andaikan $I,J$ merupakan interval dan fungsi monoton surjektif $f \colon I \to J$ memiliki invers $g \colon J \to I$.
Andaikan $x \in I$ dan $y \coloneqq f(x) \in J$, serta $g$ diferensiabel di
$y$.  Buktikan:
\begin{enumerate}[a)]
\item
Jika $g'(y) \neq 0$, maka $f$ diferensiabel di $x$.
\item
Jika $g'(y) = 0$, maka $f$ tidak diferensiabel di $x$.
\end{enumerate}
\end{exercise}
